\documentclass[12pt,a4paper]{article}
\usepackage[margin=2.5cm]{geometry}
\usepackage[T1]{fontenc}
\usepackage[utf8]{inputenc}
\usepackage[brazil]{babel}
\usepackage{lmodern}
\usepackage{hyperref}
\usepackage{enumitem}

\title{Sistema Instituto Alento \\ Manual de Uso}
\author{Instituto Alento}
\date{\today}

\begin{document}
\maketitle
\tableofcontents
\newpage

\section{Acesso ao sistema}
\subsection{Login}
\begin{enumerate}
\item Acesse \texttt{/login}.
\item Informe usuario e senha.
\item Clique em \textit{Entrar}.
\end{enumerate}

\subsection{Recuperacao de acesso}
Se o usuario nao conseguir entrar:
\begin{itemize}[leftmargin=*]
\item conferir se a conta esta ativa;
\item validar perfil e permissao com o Admin;
\item redefinir senha, quando aplicavel.
\end{itemize}

\section{Uso por perfil}
\subsection{Admin}
\begin{enumerate}
\item Abrir \texttt{/painel} para ver indicadores.
\item Entrar em \texttt{/familias} para cadastro e consulta.
\item Entrar em \texttt{/agenda} para criar ou editar compromissos.
\item Entrar em \texttt{/agenda/presencas} para registrar presenca/falta.
\item Entrar em \texttt{/administracao} para manter salas, justificativas, campos extras e filtros.
\end{enumerate}

\subsection{Medico/Voluntario}
\begin{enumerate}
\item Abrir agenda do dia.
\item Conferir pacientes e horarios.
\item Registrar presenca e observacoes.
\item Ajustar disponibilidade, quando autorizado.
\end{enumerate}

\subsection{Familia}
\begin{enumerate}
\item Consultar agenda e notificacoes.
\item Solicitar ou acompanhar consulta do dependente.
\item Atualizar dados de contato em \textit{Meus Dados}.
\end{enumerate}

\section{Passo a passo: agendar consulta}
\begin{enumerate}
\item Acessar \texttt{/agenda}.
\item Clicar em \textit{Novo Agendamento}.
\item Preencher: titulo, data, hora, tipo, profissional.
\item Selecionar sala quando o tipo de atendimento exigir.
\item Vincular familia e dependente.
\item Salvar.
\end{enumerate}

\section{Passo a passo: registrar presenca}
\begin{enumerate}
\item Acessar \texttt{/agenda/presencas}.
\item Escolher o compromisso.
\item Marcar status: presente, falta ou falta justificada.
\item Salvar observacao, se necessario.
\end{enumerate}

\section{Regras importantes de preenchimento}
\begin{itemize}[leftmargin=*]
\item Campos com asterisco sao obrigatorios.
\item Em formularios com validacao em tempo real, corrigir campos destacados antes de salvar.
\item Evitar abreviacoes confusas em nomes e observacoes.
\end{itemize}

\section{Checklist de qualidade operacional}
\begin{itemize}[leftmargin=*]
\item Todo compromisso deve ter data e hora validas.
\item Toda falta deve ter status correto e justificativa quando existir.
\item Toda alteracao sensivel deve ser feita por perfil autorizado.
\end{itemize}

\section{Erros comuns e como resolver}
\subsection{Nao consigo salvar formulario}
Revisar campos destacados em vermelho e tentar novamente.

\subsection{Nao aparece opcao de menu}
Usuario pode estar sem permissao para aquele recurso.

\subsection{Nao aparece sala livre}
Trocar horario ou confirmar se existem salas ativas no cadastro.

\end{document}
